\documentclass[12pt]{article}
\usepackage[paperheight=9in,paperwidth=9in,margin=0.5in]{geometry}
%\usepackage{fullpage}

\usepackage{mathtools}
\usepackage{times}
\usepackage{amsfonts}
\usepackage{amsmath}
\usepackage{braket}
\usepackage[x11names]{xcolor}
\usepackage{tikz}
\usepackage{titlesec}
\usepackage{background}
\usepackage{blindtext}
\backgroundsetup{ scale=1, angle=0, opacity=1, contents={\begin{tikzpicture}[remember picture,overlay] \path [left color=Orchid1,right color=white] (current page.south west)rectangle (current page.north east);  \end{tikzpicture}} }


\titleformat*{\section}{\LARGE\bfseries}
\titleformat*{\subsection}{\Large\bfseries}
\setcounter{section}{2}
%\setcounter{subsection}{1}
\pagestyle{empty}
\begin{document}
  \pagebreak 

\section{Power Mean Inequality}

\subsection{Statement :-}
\Large{\textbf{The Power Mean Inequality is a generalized form of the multi-variable QM-AM-GM-HM Inequality.}

\textbf{For }$\boldsymbol{n}$\textbf{ positive real numbers }$\boldsymbol{a_i,}$ $\boldsymbol{i=1,2,\cdots,n}$\textbf{ and }$\boldsymbol{n}$\textbf{ positive real weights }$\boldsymbol{w_i}$\textbf{ with sum }$\boldsymbol{\displaystyle{\sum\limits_{i=1}^n} w_i=1}$\textbf{, the power mean with exponent }$\boldsymbol{t}$\textbf{, where }$\boldsymbol{t\in\mathbb{R}}$\textbf{, is defined by}\[\boldsymbol{M(t)= \begin{cases} \displaystyle{\prod\limits_{i=1}^n }a_i^{w_i} &\text{if } t=0 \\ \left(\displaystyle{\sum\limits_{i=1}^n w_ia_i^t} \right)^{\dfrac{1}{t}} &\text{otherwise} \end{cases}}\]\textbf{Then the Power Mean Inequality states that for all real numbers }$\boldsymbol{k_1}$\textbf{ and }$\boldsymbol{k_2}$\textbf{, }$\boldsymbol{M(k_1)\ge M(k_2)}$\textbf{ if }$\boldsymbol{k_1>k_2}$\textbf{.}

\textbf{In particular, for non zero }$\boldsymbol{k_1}$\textbf{ and }$\boldsymbol{k_2}$\textbf{, and equal weights }$\boldsymbol{\left(\textbf{i.e. }w_i=\dfrac{1}{n}\right)}$\textbf{, if }$\boldsymbol{k_1>k_2}$\textbf{, then}\[\boldsymbol{\left( \frac{1}{n} \sum_{i=1}^n a_{i}^{k_1} \right)^{\dfrac{1}{k_1}}  \ge \left( \frac{1}{n} \sum_{i=1}^n a_{i}^{k_2} \right)^{\dfrac{1}{k_2}}}\]}\textbf{Considering the limiting behavior, we also have }$\boldsymbol{\displaystyle{\lim\limits_{t\to +\infty}} M(t)=\max\{a_i\}}$\textbf{, }\linebreak$\boldsymbol{\displaystyle{\lim\limits_{t\to -\infty}} M(t)=\min\{a_i\}}$\textbf{ and }$\boldsymbol{\displaystyle{\lim\limits_{t\to 0}}\, M(t)= M(0)}$
\end{document}


